\chapter{Introduction}
\label{ch:base-introduction}

The main manuscript develops the theory at one fixed base point \(c_*\): the design is \(\{c_*\}\),
every agent carries one common frame, and the only transports are the independent interaction links.
This supplement collects what the theory requires once contexts vary, once networks become infinite,
and once agents are placed on a lattice. It answers one question for each extension: which statements
of the single-point theory survive, and exactly which additional data (a connection, a cover, a design,
a base measure, a scale identification, or a lattice realization) each extension consumes. Base
transport is deferred until the single-point theory is complete, and nothing in the main manuscript
depends on this supplement. The relative frame \(h_i\) between belief and model frames, which the main
manuscript removes by a gauge choice because no energy, kernel, or evaluator at \(c_*\) reads it,
returns to the main manuscript only if frames become dynamical agent states.

Part B-I treats base transport and multi-point designs. \Cref{ch:base-geometry} introduces base curves,
relative frames, chosen connections and covariant defects, the \v{C}ech picture, finite designs with
their shared-link stabilizer (\Cref{prop:gen-shared-link-stabilizer}), and scale maps over base maps.
\Cref{ch:pullback-geometry} pulls the fiber Fisher and Amari--Chentsov tensors back to the base through
a chosen connection and proves the exact signed comparison of fine and coarse pullbacks
(\Cref{thm:pb-signed-base-comparison}) together with the isotropy criterion for a vanishing horizontal
defect (\Cref{thm:pb-isotropy-criterion}). \Cref{ch:section-space-histories} lifts inference histories to
spaces of sections, and \Cref{ch:neighborhood-parents} constructs parents over covers and states the
connection-descent criterion. Part B-II, beginning with \Cref{ch:renormalization} and continuing in
\Cref{ch:wilsonian-limits}, holds the asymptotic apparatus that the main manuscript does not use: mode
exponents, tempered comparisons, interaction beta functions, fixed points, spectra, and two-index
limits. Part B-III holds the lattice cluster, \Cref{ch:constructive-spatial-rg},
\Cref{ch:normalized-geometry}, and \Cref{ch:thermodynamic-fisher-draft}, whose typing, as agents at
\(c_*\) or as base points, is \Cref{open:rg-lattice-typing}. The appendices hold the proofs of the
base-direction classifications (\Cref{app:base-proofs}), of the constructive spatial family
(\Cref{app:spatial-rg-proofs}), and of the thermodynamic Fisher theorems
(\Cref{app:thermodynamic-fisher-draft-proofs}), followed by the notation table, the open problems, and
the numerical index of this supplement (\Cref{app:base-backmatter}).

The supplement cites the main manuscript and Supplement G; its chapters are numbered B0, B1, and so
on, so that references into the other two documents are unambiguous. It uses the symbols of the main
manuscript with their meanings there and adds its own. The horizontal lift \(H^x_e\), the relative frame
\(h_i\), the design index \(d\) with incidence set \(I_D\), the scale bundle morphism \(\Psi\), the
covariant defect \(\mathcal D\Phi\) of a cross map, and the base-cocycle residual \(\mathcal R\) belong
to this supplement. Here \(f\) denotes a base map and \(s\) a section; the fast and slow blocks of the
main manuscript do not occur. Frame re-expressions follow the convention of the main manuscript,
\(u'=u\cdot g^{-1}\), separately in each channel when two atlases are used, and the common-frame case is
\(g_i^b=g_i^m=g_i\). The lattice part declares its local symbols in one table at its head.

The contextual base is not identified with physical space or spacetime; graph links are not identified
with base transports except under the declared realization \Cref{hyp:geo-graph-base-transport}; no
probability law on a continuum space of sections is constructed; and the asymptotic Wilsonian vocabulary
is not applied to a finite network at \(c_*\), where finite depth has only finite-step rates
(\Cref{prop:rg-finite-depth-no-asymptotics}). \status{NOT-CLAIMED} The continuum law theory is \Cref{open:prob-continuum-theory} and the lattice typing
is \Cref{open:rg-lattice-typing}.
